%%%%%%%%%%%%%%%%%%%%%%%%%%%%%%%%%%%%%%%%%%%%%%%%%%%%%%%
%   ACL 2026 Findings Poster - UrBLiMP                %
%   Uni Konstanz colours + logo, NO themeKonstanz.sty %
%   Compile with: XeLaTeX                             %
%   Needs: graphics/UniKonstanz_Logo_Optimum.pdf      %
%          fonts: Arial (system), Noto Nastaliq Urdu  %
%%%%%%%%%%%%%%%%%%%%%%%%%%%%%%%%%%%%%%%%%%%%%%%%%%%%%%%

\documentclass[a0paper,portrait,fontscale=0.40]{baposter}

\usepackage{fontspec}
\usepackage[bidi=basic]{babel}
\babelprovide[import]{arabic}

\usepackage{amssymb}
\usepackage{booktabs}
\usepackage{array}
\usepackage{makecell}
\usepackage{graphicx}
\usepackage[table]{xcolor}
\usepackage{enumitem}
\setlist{nosep,leftmargin=*}

\usepackage{tikz}
\usetikzlibrary{positioning,shapes.geometric,backgrounds,fit,arrows.meta,calc}
\usepackage{pgfplots}
\pgfplotsset{compat=1.18}

% ---------- Uni Konstanz Corporate Design colours ----------
\definecolor{seeblau100}{RGB}{0,165,213}
\definecolor{seeblau65} {RGB}{89,194,225}
\definecolor{seeblau35} {RGB}{166,221,235}
\definecolor{seeblau20} {RGB}{205,236,243}
\definecolor{schwarz60}{gray}{0.40}
\definecolor{schwarz40}{gray}{0.60}
\definecolor{schwarz20}{gray}{0.80}
\definecolor{schwarz10}{gray}{0.90}
% -----------------------------------------------------------

\setmainfont{Arial}
\setsansfont{Arial}
\babelfont[arabic]{rm}{Noto Nastaliq Urdu}

% ---------- Phenomenon card (visual) ----------
% \phencard{<icon glyph>}{<phenomenon name>}{<good Urdu>}{<bad Urdu>}
\newcommand{\phencard}[4]{%
\begin{tikzpicture}[baseline=(card.north)]
  \node[draw=seeblau35,line width=1pt,rounded corners=4pt,
        fill=seeblau20!40,inner sep=6pt,
        text width=0.30\linewidth,align=left] (card) {%
    \begin{tikzpicture}
      \node[circle,fill=seeblau100,text=white,inner sep=3pt,
            minimum size=18pt,font=\bfseries\small] (ic) {#1};
      \node[right=4pt of ic,font=\small\bfseries,
            text=seeblau100,text width=0.78\linewidth]
            {#2};
    \end{tikzpicture}\\[2pt]
    \footnotesize
    \textcolor{seeblau100}{$\checkmark$}~\foreignlanguage{arabic}{#3}\\
    \textcolor{schwarz60}{$\times$}~\foreignlanguage{arabic}{#4}%
  };
\end{tikzpicture}%
}

\begin{document}

\begin{poster}{
  grid=false,
  columns=2,
  headerborder=closed,
  colspacing=1.2em,
  bgColorOne=white,
  borderColor=seeblau100,
  headerColorOne=seeblau100,
  headerColorTwo=seeblau65,
  headerFontColor=white,
  boxColorOne=white,
  textborder=rounded,
  eyecatcher=true,
  headerheight=0.16\textheight,
  headershape=rounded,
  headerfont=\Large\bfseries,
  boxshade=plain,
  background=none,
  linewidth=2pt
}
%%% LEFT EYE CATCHER -- Uni Konstanz logo (with safe fallback)
{\IfFileExists{graphics/UniKonstanz_Logo_Optimum.pdf}%
  {\includegraphics[height=8em]{graphics/UniKonstanz_Logo_Optimum.pdf}}%
  {\colorbox{seeblau100}{\parbox{8em}{\centering\color{white}\bfseries
    Universit\"at\\Konstanz}}}}
%%% TITLE
{\color{seeblau100}\bfseries UrBLiMP: A Benchmark for Evaluating the Linguistic
  Competence of Large Language Models in Urdu}
%%% AUTHORS (wrapped in parbox so baposter keeps the whole block together)
{\color{black}\normalsize
\parbox{0.95\linewidth}{\centering
 Farah Adeeba$^{1,2}$ \quad Brian Dillon$^{1}$ \quad Hassan Sajjad$^{3}$ \quad Rajesh Bhatt$^{1}$\\[3pt]
 \footnotesize
 $^{1}$Department of Linguistics, UMass Amherst, USA \quad
 $^{2}$Cluster of Excellence ``The Politics of Inequality'', University of Konstanz, Germany \quad
 $^{3}$Faculty of Computer Science, Dalhousie University, Halifax, Canada \\[2pt]
 \textit{ACL 2026 Findings} \;$\bullet$\; \texttt{farah.adeeba@uni-konstanz.de}}}
%%% RIGHT EYE CATCHER -- venue badge
{\color{seeblau100}\Large\bfseries ACL\\[-2pt]2026}


%====================================================================
% ROW 1
%====================================================================
\headerbox{Motivation \& Contributions}{name=motiv,column=0,row=0}{
\textcolor{seeblau100}{\textbf{Why Urdu?}}\\
LLMs are mostly evaluated on English. For low-resource, morphologically rich,
head-final languages such as \textbf{Urdu} (Indo-Aryan, $\sim$230M speakers,
Perso-Arabic script), it is unclear which grammatical phenomena multilingual
LLMs actually learn. Existing minimal-pair benchmarks (BLiMP, CLiMP, SLING,
RuBLiMP, TurBLiMP, BLiMP-NL, MultiBLiMP) lack broad coverage for Urdu.\\[4pt]
\textcolor{seeblau100}{\textbf{Our contributions}}
\begin{itemize}
\item First broad-coverage minimal-pairs benchmark for \textbf{Urdu syntax}.
\item \textbf{5{,}696} pairs, \textbf{10} phenomena, \textbf{19} paradigms.
\item Typologically distinctive coverage: split ergativity, dative
      experiencers, honorific agreement, obliqueness, participial relatives,
      word-order variation.
\item \textbf{Two-round human validation} with 17 native speakers
      (1{,}077 pairs).
\item Evaluation of \textbf{20+ LLMs}: BERT, mT5, BLOOMZ, Granite, Gemma-3,
      LLaMA-3, DeepSeek-R1, Alif-Urdu, GPT-4o.
\end{itemize}
}

\headerbox{UrBLiMP at a Glance}{name=glance,column=1,row=0}{
\centering
\begin{tikzpicture}
  \foreach \v/\l/\x in {%
     5696/{minimal pairs}/0,
     10/{phenomena}/1,
     19/{paradigms}/2,
     17/{annotators}/3,
     20/{LLMs}/4}{
    \node[circle,draw=seeblau100,line width=1.5pt,fill=seeblau20!50,
          minimum size=2.2cm,align=center] at (\x*2.4,0)
          {\textcolor{seeblau100}{\bfseries\large\v}\\[-2pt]
           \footnotesize \l};
  }
\end{tikzpicture}\\[6pt]
\raggedright\footnotesize
\textbf{Sources:} Urdu Treebank (CLE, 7{,}854 sentences) + in-house Urdu corpus.\\
\textbf{Validation:} two-round forced-choice acceptability on PCIbex.\\
\textbf{Code \& data:} \texttt{github.com/farahadeeba/UrBLiMP}
}


%====================================================================
% ROW 2 -- Linguistic phenomena AS VISUAL CARDS
%====================================================================
\headerbox{Linguistic Phenomena -- 10 cards, one per phenomenon}%
{name=phen,column=0,below=motiv,span=2}{
\centering
% Row 1 of cards
\phencard{A}{Aspect Agreement}%
  {وہ ملتا \underline{رہا} تھا}%
  {وہ ملتا \underline{چکا} تھا}\hfill
\phencard{D}{Dative Object}%
  {ماجرا راجا \underline{کو} سنایا}%
  {ماجرا راجا سنایا}\hfill
\phencard{E}{Ergativity}%
  {نے ماجرا \underline{سنایا}}%
  {نے ماجرا \underline{سناتا}}\\[6pt]
% Row 2 of cards
\phencard{X}{Experiencer Subject}%
  {\underline{مجھے} پسند آیا}%
  {\underline{میں} پسند آیا}\hfill
\phencard{H}{Honorific}%
  {بہن جی ناراض \underline{تھیں}}%
  {بہن جی ناراض \underline{تھی}}\hfill
\phencard{N}{N--J Agreement}%
  {بیس فٹ \underline{اونچا} ہے}%
  {بیس فٹ \underline{اونچی} ہے}\\[6pt]
% Row 3 of cards
\phencard{O}{Obliqueness}%
  {\underline{بڑے} صاحب نے کہا}%
  {\underline{بڑا} صاحب نے کہا}\hfill
\phencard{P}{Participial Relatives}%
  {\underline{کھولا} گیا دروازہ}%
  {\underline{کھولتا} گیا دروازہ}\hfill
\phencard{S}{Subject--Verb Agr}%
  {باپ پاگل ہو \underline{جاتا}}%
  {باپ پاگل ہو \underline{جاتی}}\\[6pt]
% Row 4 (centered single card)
\phencard{W}{Word-Order Variation}%
  {جو \underline{بات} ہے وہ بولو}%
  {جو ہے وہ بولو بات}\hfill
\begin{tikzpicture}[baseline=(legend.north)]
  \node[draw=seeblau100,line width=1pt,rounded corners=4pt,
        fill=white,inner sep=8pt,text width=0.62\linewidth] (legend) {%
    \small\textbf{How to read a card:} each badge shows the
    \textcolor{seeblau100}{$\checkmark$}~\textbf{grammatical} sentence and the
    \textcolor{schwarz60}{$\times$}~\textbf{ungrammatical} minimal-pair
    counterpart, differing by one underlined morpheme.};
\end{tikzpicture}
}


%====================================================================
% ROW 3 -- Method (visual 3-step pipeline)
%====================================================================
\headerbox{Method -- end-to-end pipeline}{name=method,column=0,below=phen,span=2}{
\centering
\begin{tikzpicture}[
   node distance=10pt and 14pt,
   stepbox/.style={draw=seeblau100,line width=1.2pt,fill=seeblau20!40,
                   rounded corners=6pt,inner sep=8pt,
                   text width=0.27\linewidth,align=left,font=\footnotesize},
   steparrow/.style={-Stealth,line width=2pt,draw=seeblau100}
]
  \node[stepbox] (s1) {%
    \textcolor{seeblau100}{\textbf{1. Source corpora}}\\[3pt]
    \textbf{Urdu Treebank} (CLE, 7{,}854 sentences) -- extract via syntactic
    patterns, e.g.\ \texttt{(VC (VBF) (AUXP))}.\\[3pt]
    \textbf{In-house corpus} -- regex extraction
    (\foreignlanguage{arabic}{جو\dots وہ\dots}) for sparse paradigms.};
  \node[stepbox,right=of s1] (s2) {%
    \textcolor{seeblau100}{\textbf{2. Pair generation}}\\[3pt]
    Rule-based perturbation of \emph{one} minimal feature
    (swap aspect aux., drop \textit{ko}, change verb gender, \dots) to derive
    the \textbf{ungrammatical} counterpart.\\[3pt]
    Manual review + typo cleaning per paradigm.};
  \node[stepbox,right=of s2] (s3) {%
    \textcolor{seeblau100}{\textbf{3. Human validation (PCIbex)}}\\[3pt]
    \textbf{17 native Urdu speakers}, two rounds.\\
    R1: 20-pair training set.\\
    R2: $\sim$190 pairs/annotator, 10/paradigm, randomised, forced-choice.\\[3pt]
    Validated subset: \textbf{1{,}077 pairs}.};
  \draw[steparrow] (s1) -- (s2);
  \draw[steparrow] (s2) -- (s3);
\end{tikzpicture}
}


%====================================================================
% ROW 4 -- Evaluation setup + headline results
%====================================================================
\headerbox{Evaluation Setup}{name=eval,column=0,below=method}{
\textcolor{seeblau100}{\textbf{Models (20+)}}
\begin{itemize}
\item \textbf{Encoder:} mBERT
\item \textbf{Encoder--decoder:} mT5-small, mT5-large
\item \textbf{Decoder:} Gemma-3 (1B/4B/12B/27B, pt+it), LLaMA-3-8B/70B,
      BLOOMZ-7B, Granite-3.3 (2B/8B), DeepSeek-R1-Distill-LLaMA-70B
\item \textbf{Urdu-tuned:} Alif-1.0-8B-Instruct
\item \textbf{Closed:} GPT-4o (Base, Grammar-Prompt, GP+CoT)
\end{itemize}
\textcolor{seeblau100}{\textbf{Metric}} -- sentence-level accuracy:
\[\text{Acc} = \tfrac{1}{N}\sum_{i=1}^{N}\mathbb{I}\!\left[
   \text{ppl}(s^{i}_{\text{good}}) < \text{ppl}(s^{i}_{\text{bad}})\right]\]
GPT-4o evaluated via two-choice prompting (\textit{explain-then-process}).
}

\headerbox{Headline Results -- Avg.\ Accuracy across 10 Phenomena}%
{name=results,column=1,below=method}{
\centering
\begin{tikzpicture}
\begin{axis}[
   xbar,
   width=0.95\linewidth,height=12cm,
   bar width=8pt,
   enlarge y limits=0.04,
   xmin=45,xmax=100,
   xlabel={\footnotesize Average accuracy (\%)},
   xtick={50,60,70,80,90,100},
   ytick=data,
   symbolic y coords={mT5-small,mT5-large,mBERT,Granite-8B,LLaMA-3-8B,
        DeepSeek-70B,Gemma-3-12B,Gemma-3-27B,Alif-8B,GPT-4o (CoT),
        LLaMA-3-70B,GPT-4o (Base),Human,GPT-4o (GP)},
   nodes near coords,nodes near coords align={horizontal},
   nodes near coords style={font=\scriptsize},
   tick label style={font=\scriptsize},
   label style={font=\footnotesize},
]
\addplot[fill=seeblau100,draw=seeblau100] coordinates {
   (50.69,mT5-small) (63.36,mT5-large) (78.77,mBERT)
   (87.44,Granite-8B) (90.71,LLaMA-3-8B) (92.46,DeepSeek-70B)
   (92.68,Gemma-3-12B) (92.42,Gemma-3-27B) (93.25,Alif-8B)
   (94.25,GPT-4o (CoT)) (94.73,LLaMA-3-70B) (95.69,GPT-4o (Base))
   (97.40,GPT-4o (GP))};
\addplot[fill=seeblau35,draw=seeblau100] coordinates {
   (96.28,Human)};
\draw[red,dashed,thick] (axis cs:96.28,mT5-small)
   -- (axis cs:96.28,GPT-4o (GP)) node[pos=0.9,above,
   font=\scriptsize,red] {Human baseline};
\end{axis}
\end{tikzpicture}\\[-4pt]
{\footnotesize\textit{Only \textbf{GPT-4o + Grammar-Prompt} exceeds the
human baseline. LLaMA-3-70B \& Urdu-tuned Alif-8B are within reach.}}
}


%====================================================================
% ROW 5 -- Findings: model size chart + 2 text findings
%====================================================================
\headerbox{Key Findings}{name=findings,column=0,below=eval,span=2}{
\begin{minipage}[t]{0.36\linewidth}
\centering
\textcolor{seeblau100}{\textbf{Size $\neq$ better}}\\[4pt]
\begin{tikzpicture}
\begin{axis}[
   ybar,
   width=\linewidth,height=5.5cm,
   bar width=14pt,
   ymin=80,ymax=100,
   ylabel={\footnotesize Acc.\ (\%)},
   symbolic x coords={1B,4B,12B,27B},
   xtick=data,
   nodes near coords,
   nodes near coords style={font=\scriptsize},
   tick label style={font=\scriptsize},
   label style={font=\footnotesize},
   ymajorgrids=true,grid style=dashed,
]
\addplot[fill=seeblau100,draw=seeblau100] coordinates {
   (1B,86.74) (4B,88.11) (12B,92.68) (27B,92.42)};
\draw[red,dashed,thick] (axis cs:1B,96.28) -- (axis cs:27B,96.28)
   node[pos=0.5,above,font=\scriptsize,red] {Human};
\end{axis}
\end{tikzpicture}\\[-4pt]
{\footnotesize Gemma-3-PT 1B$\to$27B: significant gains
(Wilcoxon $p<0.05$). But 12B $\approx$ 27B $\approx$ LLaMA-70B
$\approx$ Alif-8B (no sig.\ diff.\ among top models).}
\end{minipage}\hfill
\begin{minipage}[t]{0.30\linewidth}
\textcolor{seeblau100}{\textbf{Long-distance agreement is hard}}\\[3pt]
Even LLaMA-3-70B fails when subject \& verb are separated by intervening
clauses:\\[2pt]
\foreignlanguage{arabic}{استانی سبق پڑھا کر \dots \underline{سوتی}}\\[2pt]
{\footnotesize The longer the dependency, the more LMs default to local
cues.}\\[6pt]
\textcolor{seeblau100}{\textbf{Pretrained $>$ Instruction-tuned}}\\[3pt]
On Gemma-3 (1B/4B/12B): instruction tuning can \emph{hurt} grammatical
sensitivity in 8/10 phenomena.
\end{minipage}\hfill
\begin{minipage}[t]{0.30\linewidth}
\textcolor{seeblau100}{\textbf{Phenomenon-specific}}
\begin{itemize}
\item \textbf{Aspect, Dative:} near-perfect (local cues).
\item \textbf{Obliqueness:} hard for adjectives \& masc.\ sg.\ nouns
      (subtle morphology).
\item \textbf{Subj--Verb gender:} hardest, max 92.5\%.
\item \textbf{Honorific:} requires register knowledge.
\item \textbf{UrBLiMP-SVA} is harder than MultiBLiMP-SVA for every model.
\end{itemize}
\end{minipage}
}


%====================================================================
% ROW 6 -- Take-aways
%====================================================================
\headerbox{Take-aways}{name=take,column=0,below=findings,span=2}{
\begin{itemize}
\item UrBLiMP is the first broad-coverage minimal-pair benchmark for
      \textbf{Urdu syntax}.
\item Multilingual LLMs handle \textbf{local agreement} well, but
      struggle with \textbf{long-distance dependencies} and
      \textbf{morpho-syntactic case}.
\item \textbf{Grammar prompting} brings GPT-4o above the human average.
\item \textbf{Continual pretraining on Urdu} (Alif-8B) beats its
      LLaMA-3-8B base on 9/10 phenomena.
\item Resources: \texttt{github.com/farahadeeba/UrBLiMP} -- open-source.
\end{itemize}
\vfill
\centering
{\footnotesize\textcolor{seeblau100}{\textbf{\textit{ACL 2026 Findings
   \;$\bullet$\; uni-konstanz.de \;$\bullet$\; umass.edu \;$\bullet$\; dal.ca}}}}
}

\end{poster}
\end{document}
